\section{The ball Dirichlet problem with continuous density}
\label{sec:bt-continuous-ball-dirichlet}

The linear density inequalities make subsolutions stable under maxima and controlled quadratic perturbations. These operations allow the homogeneous ball envelope to be replaced by an envelope with prescribed density. The exact-center automorphisms still give bounded second differences, provided their Jacobian factor is included in the density.

\begin{lemma}\label{lem:continuous-density-subsolution-algebra}
Let $u,v\in\PSH(\Omega)\cap L^\infty_{\loc}$ and let $F,G\ge0$ be continuous with $(\ddc u)^n\ge F^n\eta^n$ and $(\ddc v)^n\ge G^n\eta^n$. For $0\le t\le1$ and $a\ge0$,
\[
 \bigl(\ddc((1-t)u+tv+a|z|^2)\bigr)^n
 \ge\bigl((1-t)F+tG+a\bigr)^n\eta^n.
\]
Also $(\ddc\max\{u,v\})^n\ge\min\{F,G\}^n\eta^n$.
\end{lemma}
\begin{proof}
For every admissible $B$, the linear inequalities in \cref{prop:bounded-monge-ampere-linear-density} give
\[
 L_B((1-t)u+tv+a|z|^2)
 \ge n((1-t)F+tG)+a\operatorname{tr}B
 \ge n((1-t)F+tG+a).
\]
Here $\operatorname{tr}B\ge n$ by arithmetic--geometric mean. The same proposition gives the first assertion. An upper test for $\max\{u,v\}$ at a point touches at least one of $u,v$ there. Its determinant is at least the corresponding density, and hence at least $\min\{F,G\}^n$. Apply the upper-test characterization again.
\end{proof}

\begin{theorem}\label{thm:continuous-density-ball-dirichlet}
Let $B=B(a,R)\subset\mathbb C^n$, $n\ge1$. For $f\in C(\overline B)$ with $f\ge0$ and real $\varphi\in C(\partial B)$, there is a unique $H_{f,\varphi}\in\PSH(B)\cap C(\overline B)$ with
\[
 (\ddc H_{f,\varphi})^n=f\eta^n,
 \qquad H_{f,\varphi}|_{\partial B}=\varphi.
\]
It is the pointwise supremum of bounded plurisubharmonic $v$ satisfying $(\ddc v)^n\ge f\eta^n$ and upper boundary limits at most $\varphi$. For any other such data $g,\chi$,
\begin{equation}\label{eq:ball-dirichlet-root-stability}
 \|H_{f,\varphi}-H_{g,\chi}\|_{L^\infty(B)}
 \le\|\varphi-\chi\|_{L^\infty(\partial B)}
       +R^2\|f^{1/n}-g^{1/n}\|_{L^\infty(B)}.
\end{equation}
If $f^{1/n}$ and $\varphi$ have real $C^2$ extensions to neighborhoods of $\overline B$ and $\partial B$, respectively, then $H_{f,\varphi}\in C^{1,1}_{\loc}(B)$.
\end{theorem}
\begin{proof}
An affine change $z=a+R\zeta$ reduces the problem to the unit ball. The transformed density is $R^{2n}f(a+R\zeta)$, so its root is $R^2f^{1/n}(a+R\zeta)$. We prove the assertions there and put $\psi=|z|^2-1$ and $F=f^{1/n}$.

\emph{Step 1. The subsolution envelope for regular data.}
First suppose $F$ and $\varphi$ have the stated $C^2$ extensions. Extend $\varphi$ as a real $C^2$ function $b$ on a neighborhood of $\overline B$. For a sufficiently large $A$, $s=b+A\psi$ is strictly plurisubharmonic and $\det H(s)\ge f$ on $B$. It is a bounded subsolution with the required boundary value. Every bounded competitor is at most the harmonic extension $P\varphi$, by \cref{thm:subharmonic-comparison}. Let $H$ be their supremum. By \cref{lem:continuous-density-subsolution-algebra}, finite maxima of competitors, including $s$, remain competitors. Choquet selection in \cref{thm:psh-upper-envelopes} supplies such finite maxima increasing almost everywhere to the upper regularization $H^*$. They lie between $s$ and $P\varphi$. Increasing continuity in \cref{thm:bt-increasing-continuity} therefore gives $(\ddc H^*)^n\ge f\eta^n$. The same two barriers give boundary value $\varphi$, so $H^*$ is a competitor. Hence $H=H^*$ pointwise, and $H$ is bounded and plurisubharmonic.

\emph{Step 2. Second differences and equality of the density.}
Fix a compact set of centers $K\Subset B$ and use $S_{a,h}$ from \cref{lem:ball-centered-automorphisms}. Write $J_{a,h}(z)=\det_{\mathbb C}D_zS_{a,h}$. Local smooth approximation and decreasing continuity show that Bedford--Taylor measures commute with these biholomorphisms. Consequently
\[
 \bigl(\ddc(H\circ S_{a,h})\bigr)^n
 \ge (F\circ S_{a,h})^n|J_{a,h}|^2\eta^n.
\]
The explicit formula for the automorphisms has denominators bounded away from zero on $\overline B$, uniformly for $a\in K$ and sufficiently small $h$. Its inverse has the same property. Thus $|J_{a,h}|^{2/n}$ is smooth and strictly positive there, with its first two real $h$ derivatives uniformly bounded. Taylor's formula and $S_{a,0}=\mathrm{id}$ give constants $C_1,C_2\ge0$ such that
\[
 \begin{aligned}
 \tfrac12\bigl((F\circ S_{a,h})|J_{a,h}|^{2/n}
 +(F\circ S_{a,-h})|J_{a,-h}|^{2/n}\bigr)
 &\ge F-C_1|h|^2 &&\text{on }\overline B,\\
 \tfrac12(\varphi\circ S_{a,h}+\varphi\circ S_{a,-h})
 &\le\varphi+C_2|h|^2 &&\text{on }\partial B.
 \end{aligned}
\]
The constants depend on $K,n$ and the fixed $C^2$ data, independently of $a,h,z$. By \cref{lem:continuous-density-subsolution-algebra},
\[
 V_{a,h}=\tfrac12(H\circ S_{a,h}+H\circ S_{a,-h})
            +C_1|h|^2\psi-C_2|h|^2
\]
has Monge--Amp\`ere measure at least $f\eta^n$ and upper boundary limits at most $\varphi$. Indeed the average's root lower bound is the nonnegative average in the first line above; adding $C_1|h|^2\psi$ raises it by $C_1|h|^2$. Therefore $V_{a,h}\le H$. Evaluating at $a$ and using $-1\le\psi(a)\le0$ gives
\[
 H(a+h)+H(a-h)-2H(a)\le2(C_1+C_2)|h|^2.
\]
As in the proof of \cref{thm:homogeneous-ball-dirichlet}, convolution turns this into an upper bound for every real directional Hessian; plurisubharmonicity in the direction $\mathrm ih$ gives its lower bound. \Cref{prop:bounded-hessian-taylor} then gives the canonical locally $C^{1,1}$ representative and its almost-everywhere second-order Taylor expansion. Bounded Hessian coefficients, Lebesgue differentiation and uniform Bedford--Taylor continuity identify $(\ddc H)^n=\det(H_{j\bar k})\eta^n$.

At a Taylor point $x$ where $\det(H_{j\bar k})(x)>f(x)$, the complex Hessian is positive definite. Let $Q$ be the real quadratic Taylor polynomial. Choose $\varepsilon>0$ so small that $Q-\varepsilon|z-x|^2$ remains strictly plurisubharmonic and its determinant exceeds $f$ on a small ball about $x$. Shrink its radius $r$ so that the Taylor remainder on the sphere is less than $\varepsilon r^2/2$. Then
$q=Q-\varepsilon|z-x|^2+\varepsilon r^2/4$ is below $H$ on that sphere and above it at $x$. Paste $\max\{H,q\}$ inside the small ball and $H$ outside. Gluing and \cref{lem:continuous-density-subsolution-algebra} give a bounded competitor exceeding $H$, a contradiction. Since $\det(H_{j\bar k})\ge f$ almost everywhere already, equality holds almost everywhere. This proves the equation for regular data.

\emph{Step 3. Root-density stability and continuous data.}
For two regular data sets, put $\delta=\|\varphi-\chi\|_\infty$ and $\varepsilon=\|f^{1/n}-g^{1/n}\|_\infty$. The function $H_{g,\chi}+\varepsilon\psi-\delta$ has upper boundary limits at most $\varphi$ and, by \cref{lem:continuous-density-subsolution-algebra}, measure at least $(g^{1/n}+\varepsilon)^n\eta^n\ge f\eta^n$. It is a competitor for $H_{f,\varphi}$. Reverse the data sets to obtain \eqref{eq:ball-dirichlet-root-stability} with $R=1$.

For continuous data, extend $F\ge0$ continuously outside $\overline B$ by radial projection and convolve with nonnegative kernels. This gives smooth nonnegative $F_j\to F$ uniformly on $\overline B$. Smooth boundary approximations $\varphi_j\to\varphi$ are obtained on an annular neighborhood. The estimate just proved makes $H_{F_j^n,\varphi_j}$ uniformly Cauchy on $\overline B$. Its limit $U$ is continuous and plurisubharmonic, with boundary value $\varphi$. Uniform continuity of the bounded products gives $(\ddc U)^n=f\eta^n$. Applying the estimate to approximations of two sets of data and taking their uniform limits proves the stated stability bound for their solutions.

\emph{Step 4. Uniqueness and the envelope.}
Let $U$ be the solution constructed above and let $v$ be any bounded plurisubharmonic competitor for these continuous data. For $\varepsilon>0$, put $v_\varepsilon=v+\varepsilon\psi$. Its measure is at least $(F+\varepsilon)^n\eta^n$, whereas that of $U$ is $F^n\eta^n$. Boundary comparison gives
\[
 \int_{\{U<v_\varepsilon\}}(F+\varepsilon)^n\eta^n
 \le\int_{\{U<v_\varepsilon\}}F^n\eta^n.
\]
The strict difference makes this Borel set Lebesgue null. Uniqueness of canonical plurisubharmonic representatives gives $U\ge v_\varepsilon$ pointwise, as in \cref{thm:bt-comparison}. Let $\varepsilon\downarrow0$ to get $U\ge v$. Since $U$ itself is a competitor, it equals their pointwise envelope. Any other continuous solution is a competitor and is therefore at most $U$; applying the same argument with the solutions interchanged gives equality. Scaling back proves all assertions.
\end{proof}
\index{Dirichlet problem!continuous Monge--Amp\`ere density}
